\documentclass[10pt,a4paper]{amsart}
\usepackage[margin=19mm]{geometry}
\usepackage{amsmath,amssymb,mathtools}
\usepackage[scr=rsfs,cal=euler]{mathalfa}
\usepackage{xfrac}
\usepackage[unicode,hidelinks,bookmarks=false]{hyperref}
\newcommand{\ie}{i.e.\ }
\newcommand{\cf}{cf.\ }
\newcommand{\resp}{respectively\ }
\newcommand{\Subsection}{\subsection}
\providecommand{\nobreakdash}{\nobreak\hbox{-}}
\setlength{\emergencystretch}{2em}
\multlinegap0pt
\usepackage{tcolorbox}
\usepackage{amssymb}
\usepackage{xcolor}
\usepackage{mathtools}
\usepackage[shortlabels]{enumitem}
\newtheorem{lemma}{Lemma}
\newcommand{\R}{{\mathbb R}}
\newcommand{\dd}{{\mathrm d}}
\title{On high dimensional maximal functions associated to Gaussians, balls, and spheres}
\author{Valentina Ciccone, B{\l}a{\.z}ej Wr{\'o}bel}
\date{Source: arXiv:2509.13791v2 (CC BY 4.0)}
\begin{document}
\maketitle

\noindent\emph{Source and licence.} On high dimensional maximal functions associated to Gaussians, balls, and spheres. Version arXiv:2509.13791v2 (CC BY 4.0), distributed by arXiv under the Creative Commons Attribution 4.0 International licence, http://creativecommons.org/licenses/by/4.0/, as recorded in the arXiv licence metadata for this exact version and shown on the abstract page https://arxiv.org/abs/2509.13791v2. Full licence notice: \texttt{licenses/arxiv-2509.13791-CC-BY-4.0.txt}.\par\smallskip

\noindent\textit{Editorial scope.} Counted unit: the article's lemma identity\_for\_mu,g (source0.tex lines 709--716). The article's complete original proof of it is delivered with it (source0.tex lines 717--728, one \texttt{proof} environment of 12 lines). No mathematics is written, repaired or re-derived: the statement and the proof are the article's own lines, in the article's own notation, and no retained line is substituted or re-flowed.  No further source-local context is needed: every cross-reference of every retained line resolves inside the two delivered spans.  Nothing was omitted from any retained span, so the package carries no omission record.  No retained line contains a citation, so the wrapper carries no bibliography and no bibliography entry.  No retained line contains a figure environment, an image-inclusion command or drawing code, so no image is delivered and the package declares no asset dependency.  Editorial formatting only, in addition: an AMS \texttt{amsart} preamble that restates the article's own declarations and macros used by the retained lines, namely the environments \texttt{lemma} on separate counters, together with the standard packages those lines need.  The retained environments print in this document as Lemma 1 in order of appearance, whereas the article prints the same items with the numbering of its own full document.  The complete native source of the article as distributed in the arXiv e-print for this exact version is bundled unchanged as \texttt{sources/185190/source0.tex}.
\begin{lemma}\label{lem: identity_for_mu,g} The following pointwise identities hold for $\xi \in \R^d$
\begin{align}
\label{lemma_identity_for_g}
g(\xi)&=\mathop{\mathbb{E}}_{X\sim \mathcal{N}(0,\textrm{I}_d)} e^{- 2\pi i \tfrac{X}{\sqrt{d}}\cdot \xi},\\
\label{lemma_identity_for_mu}
   \mu(\xi) &= \mathop{\mathbb{E}}_{X\sim \mathcal{N}(0,\textrm{I}_d)} e^{-2\pi i \tfrac{X}{|X|}\cdot \xi}~.
\end{align}
\end{lemma}

\begin{proof} 
First we justify \eqref{lemma_identity_for_g}. Using the fact that the Fourier transform of a Gaussian function is itself a Gaussian function we have that
$$\mathop{\mathbb{E}}_{X\sim \mathcal{N}(0,\textrm{I}_d)} e^{- 2\pi i \tfrac{X}{\sqrt{d}}\cdot \xi}= \frac{1}{(2\pi)^{d/2}} \int_{\mathbb{R}^d} e^{-2\pi i \tfrac{x}{\sqrt{d}}\cdot\xi}e^{-|x|^2/2}dx = e^{-\tfrac{2\pi^2|\xi|^2}{d}}= g(\xi)~.$$

To prove \eqref{lemma_identity_for_mu} we compute
\begin{align*} \mathop{\mathbb{E}}_{X\sim \mathcal{N}(0,\textrm{I}_d)} e^{-2\pi i \tfrac{X}{|X|}\cdot \xi} & = \frac{1}{(2\pi)^{d/2}}\int_{\mathbb{R}^d}  e^{-2\pi i \tfrac{x}{|x|}\cdot \xi} e^{-\tfrac{|x|^2}{2}}\dd x \\
& = \frac{1}{(2\pi)^{d/2}}\int_0^\infty \bigg( \int_{\mathbb{S}^{d-1}}  e^{-2\pi i\omega \cdot \xi} \dd\sigma(\omega)\bigg) e^{-\tfrac{r^2}{2}} r^{d-1}\dd r  \\
&= \sigma(\mathbb{S}^{d-1}) \mu(\xi) \frac{1}{(2\pi)^{d/2}}\int_0^\infty e^{-\tfrac{r^2}{2}} r^{d-1}\dd r \\
&= \mu(\xi) \frac{1}{(2\pi)^{d/2}}\int_{\mathbb{R}^d} e^{-\tfrac{|x|^2}{2}}\dd x \\
&= \mu(\xi)~.
\end{align*}
\end{proof}

\end{document}
